\documentclass[12pt]{article}
\usepackage[utf8]{inputenc}

\usepackage{_mystyle}

\title{Oral Exam}
\author{Joseph Sullivan}
\date{September 2025}

\DeclareMathOperator{\Kom}{Kom}
\DeclareMathOperator{\Cyl}{Cyl}
\DeclareMathOperator{\Stab}{Stab}
\DeclareMathOperator{\Slice}{Slice}
\DeclareMathOperator{\chern}{ch}
\DeclareMathOperator{\todd}{td}

\DeclareMathOperator{\ext}{ext}
%\DeclareMathOperator{\hom}{hom}

\newcommand{\cQ}{\mathcal{Q}}
\newcommand{\cZ}{\mathcal{Z}}
\newcommand{\LCom}{\mathcal{LC}}

\newcommand{\rddots}{\text{\reflectbox{$\ddots$}}}


\newcommand{\fg}{\mathfrak{g}}
\newcommand{\fh}{\mathfrak{h}}
\newcommand{\gl}{\operatorname{\mathfrak{gl}}}
\newcommand{\rad}{\mathrm{rad}}

\newcommand{\fa}{\mathfrak{a}}
\newcommand{\fb}{\mathfrak{b}}
\newcommand{\fn}{\mathfrak{n}}
\newcommand{\fs}{\mathfrak{s}}

\DeclareMathOperator{\ad}{ad}
\DeclareMathOperator{\rank}{rank}

\newcommand{\so}{\mathfrak{so}}
\newcommand{\SO}{\mathrm{SO}}
\newcommand{\fsl}{\operatorname{\mathfrak{sl}}}
\DeclareMathOperator{\SL}{SL}
\DeclareMathOperator{\Fl}{Fl}
\DeclareMathOperator{\tr}{tr}

\newcommand{\fW}{\mathfrak{W}}
\newcommand{\cW}{\mathcal{W}}

\DeclareMathOperator{\Gr}{Gr}


\begin{document}

\maketitle

\section{History}

Over my time at Utah so far, the following topics are the most relevant subjects I have studied. In chronological order:
\begin{itemize}
    \item Macaulay2 and Resolving Singularities on Surfaces (following Karl's Macaulay2 seminar)
    \item Complex Geometry (following Harold's class, from Huybrechts's \textit{Complex Geometry} Chapters 1-5)
    \item Syzygies (Ein-Lazarsfeld \textit{Lectures on the Syzygies and Geometry of Algebraic Varieties} 1-6), in particular:
    \begin{itemize}
        \item Minimal Resolutions
        \item Castelnuovo-Mumford Regularity
    \end{itemize}
    \item Derived/Triangulated Categories (in algebraic geometry or otherwise)
    \begin{itemize}
        \item Definitions (Gelfand Manin \textit{Methods in Homological Algebra} II-IV)
        \item Exceptional Collections/Tilting Bundles (Calabrese's \textit{On a Theorem by Beilinson}, Craw's \textit{Explicit Methods for Derived Categories of Sheaves} Lectures 1-2)
        \item Quiver Representations (Assem, Simson, Skowronski's \textit{Elements of the Representation Theory of Associative Algebras} Chapters I-III)
        \item In Algebraic Geometry (Huybrechts \textit{FM Transforms} 1-8)
        \item t-Structures and Bridgeland Stability Conditions (Bayer's \textit{A tour to stability conditions on derived categories})
    \end{itemize}
    \item Schemes (following Harold's class, Hartshorne II.1-II.8, IV.1)
    \item Representation Theory of Complex Lie Algebras (from Dragan's class, and Fulton and Harris \textit{Representation Theory} Chapters 14, 15, 18)
    \item Cohomology Base Change (Vakil Chapters 24-25)
    \item Construction of the Hilbert Scheme (Alper's \textit{Stacks and Moduli} Chapter 1, Mumford's \textit{Lectures on Curves and Surfaces} for Flattening Stratification)
    \item Theorem on Formal Functions (Vakil Chapter 28) and applications to Derived Categories (Bondal's \textit{Projective Bundles, Monoidal Transformations, and Derived Categories of Coherent Sheaves})
\end{itemize}


\section{Interests}

I have particularly enjoyed understanding the derived categories of Fano varieties through exceptional collections/semiorthogonal decompositions. I spent a lot of time wrestling with the derived equivalence of $\Coh(\P^n)$ and representations of ``the'' Beilinson quiver (you get different options for the two different Beilinson collections, or any mutations).

I then moved on to study Homological Projective Duality, in particular to understand the relation between pencils of quadrics and hyperelliptic curves from a derived category perspective. To start, I needed to understand the derived category of just one quadric by Kapranov, which took me on a journey through Koszul rings (to generalize Beilinson's resolution of the diagonal) and representation theory (to understand the cohomology of ``spinor bundles'' by Borel-Weil-Bott). There are still many gaps in my understanding of this story, but I am thoroughly enjoying it.

I have also started trying to think about moduli. I am starting out with the Hilbert scheme, but I have my sights set also on moduli spaces of semistable complexes (with a fixed Chern character) with respect to a given Bridgeland stability condition, and also moduli spaces appearing in Quantum cohomology.

\section{Syllabus}

I hope to spend a lot of my oral exam discussing my most recent focus on the construction of the Hilbert scheme and necessary cohomology results (like ``Cohomology and Base Change'' theorems) to describe the morphism to the appropriate Grassmannian. I am also particularly fond of Beilinson's resolution of the diagonal.

I am then happy to talk about foundational algebraic geometry, derived categories, or representation theory---and then to continue wherever our  discussion goes.

\end{document}
